% SGA 3, Expose III, Loop 1 English reconstruction.
% Scope: 4.11 through Definition 4.16.1, stopping immediately before
% Lemma 4.17.
% Authority: Polo--Gille Exp3-14oct24.pdf, component-local pp. 55--62
% (combined-reader pp. 165--172; original printed pp. 136--142), with only
% the section-4.11 portion of local p. 55 admitted. Force-OCR is a locator
% only. English wording was compared with Jacob C. Reinhold's CC BY 4.0
% Markdown lineage at commit e7a259f3f8608ad3edf9bf6eead3fd504dd2d23e.
% The Polo--Gille PDF controls prose, mathematics, numbering, notes, and the
% thirteen source-derived Loop-1 diagram images included below.

\setcounter{footnote}{98}

% Source page: Exp3 local p. 55; combined p. 165; original printed p. 136.
\medskip
\noindent\SGAInlineRefTarget{sga3:iii:paragraph:4-11}\textbf{4.11.}
现在研究如下情形. 设 \(S\), \(S_{\mathscr J}\) 和 \(S_0\)
如\SGARefLink{sga3:iii:proposition:4-8}{4.8} 所述.
给定三个 \(S\)-概形 \(X,X',T\), \(X\) 的子概形 \(Y\)
(相应地, \(X'\) 的子概形 \(Y'\)), 以及态射
\(f:T\to X'\) 和 \(g:X'\to X\).

% Raster stand-in for Loop 1; replace by native TeX/TikZ in Loop 2.
% Authority: Exp3-14oct24.pdf local p. 55; combined reader p. 165.
\begin{center}
  \SGAThreeNativeInput[0.62\textwidth]
    {figures/exp3/diagram_043_exp3_p055_411_setup.tex}
\end{center}

% Source page: Exp3 local p. 56; combined p. 166; original printed p. 137.
假设模 \(\mathscr J\) 约化后, 此图可以补成交换图

% Raster stand-in for Loop 1; replace by native TeX/TikZ in Loop 2.
% Authority: Exp3-14oct24.pdf local p. 56; combined reader p. 166.
\begin{center}
  \SGAThreeNativeInput[0.72\textwidth]
    {figures/exp3/diagram_044_exp3_p056_411_reduction.tex}
\end{center}

因此, 由\SGARefLink{sga3:iii:proposition:4-8}{4.8}, 有障碍
\begin{align*}
 c(X,Y,g\circ i')&\in
 \operatorname{Hom}_{\mathscr O_{Y'_0}}
 \bigl(i_0'^*g_0^*(\mathscr N_{Y/X}
   \otimes_{\mathscr O_Y}\mathscr O_{Y_0}),
   \mathscr J\mathscr O_{Y'}\bigr),\\
 c(X',Y',f)&\in
 \operatorname{Hom}_{\mathscr O_{T_0}}
 \bigl(f_0^*(\mathscr N_{Y'/X'}
   \otimes_{\mathscr O_{Y'}}\mathscr O_{Y'_0}),
   \mathscr J\mathscr O_T\bigr),\\
 c(X,Y,g\circ f)&\in
 \operatorname{Hom}_{\mathscr O_{T_0}}
 \bigl(f_0^*g_0^*(\mathscr N_{Y/X}
   \otimes_{\mathscr O_Y}\mathscr O_{Y_0}),
   \mathscr J\mathscr O_T\bigr),
\end{align*}
我们要计算它们之间的关系.
\footnote{编者注: 从\SGARefLink{sga3:iii:lemma:4-17}{4.17} 起, 将此结果应用于如下情形:
\(X\) 为 \(S\)-群, \(g:X\times_S X\to X\) 为乘法,
\(Y\) 为 \(X\) 的子概形, 使得 \(Y_{\mathscr J}\) 是
\(X_{\mathscr J}\) 的子群, \(Y'=Y\times_S Y\),
而两个态射 \(Y^3\to X^2\) 分别将 \((y_1,y_2,y_3)\) 映到
\((y_1y_2,y_3)\) 及 \((y_1,y_2y_3)\).
此时比较上述障碍将表明, \(Y\) 成为 \(X\) 的子群的障碍属于某个
Hochschild 上同调群 \(H^2(Y_0,N_0)\).}

\medskip
\noindent\SGAInlineRefTarget{sga3:iii:lemma:4-12}\textbf{引理 4.12.}
\textit{假设 \(Y'\) 在 \(S\) 上平坦, 因而
\[
 \mathscr J\otimes_{\mathscr O_{S_0}}\mathscr O_{Y'_0}
 =\mathscr J\mathscr O_{Y'}.
\]
则:}
\begin{enumerate}
\item[(i)] \textit{存在自然态射}
\begin{align*}
 b_{f_0}:&\
 \operatorname{Hom}_{\mathscr O_{Y'_0}}
 \bigl(i_0'^*g_0^*(\mathscr N_{Y/X}
   \otimes_{\mathscr O_Y}\mathscr O_{Y_0}),
   \mathscr J\mathscr O_{Y'}\bigr)\\
 &\longrightarrow
 \operatorname{Hom}_{\mathscr O_{T_0}}
 \bigl(f_0^*g_0^*(\mathscr N_{Y/X}
   \otimes_{\mathscr O_Y}\mathscr O_{Y_0}),
   \mathscr J\mathscr O_T\bigr).
\end{align*}
\item[(ii)] \textit{还存在关于 \(T\) 具有函子性的自然态射,}
\footnote{编者注: 已删除 \(Y'_0\) 在 \(X'_0\) 中为局部完全交的假设,
由\SGARefLink{sga3:iii:remark:4-8-0}{4.8.0} 可知它是多余的.
还补充说明 \(a_{g_0}(f_0)\) \lq\lq 关于 \(T\) 具有函子性\rq\rq,
这一事实在\SGARefLink{sga3:iii:lemma:4-17}{4.17} 的证明中起关键作用.}
\begin{align*}
 a_{g_0}(f_0):&\
 \operatorname{Hom}_{\mathscr O_{T_0}}
 \bigl(f_0^*(\mathscr N_{Y'/X'}
   \otimes_{\mathscr O_{Y'}}\mathscr O_{Y'_0}),
   \mathscr J\mathscr O_T\bigr)\\
 &\longrightarrow
 \operatorname{Hom}_{\mathscr O_{T_0}}
 \bigl(f_0^*g_0^*(\mathscr N_{Y/X}
   \otimes_{\mathscr O_Y}\mathscr O_{Y_0}),
   \mathscr J\mathscr O_T\bigr).
\end{align*}
\end{enumerate}

\noindent\textit{证明.}
\footnote{编者注: 此证明作了扩充, 以明确说明 \(a_{g_0}\)
\lq\lq 关于 \(T\) 的函子性\rq\rq.}
首先注意, 固定 \(X,X',Y,Y'\) 后, 给定上述 \(T\)
等价于给定态射 \((f,f_{\mathscr J}):T\to X'\times_{X'^{+}}Y'^{+}\).
令 \(Z=X'\times_{X'^{+}}Y'^{+}\), 并设 \(M\) 和 \(M'\) 为
如下定义的 \(Z\)-函子: 对每个 \(f:T\to Z\), 令
\begin{align*}
 M(T)&=
 \operatorname{Hom}_{\mathscr O_{T_0}}
 \bigl(f_0^*g_0^*(\mathscr N_{Y/X}
   \otimes_{\mathscr O_Y}\mathscr O_{Y_0}),
   \mathscr J\mathscr O_T\bigr),\\
 M'(T)&=
 \operatorname{Hom}_{\mathscr O_{T_0}}
 \bigl(f_0^*(\mathscr N_{Y'/X'}
   \otimes_{\mathscr O_{Y'}}\mathscr O_{Y'_0}),
   \mathscr J\mathscr O_T\bigr).
\end{align*}
\footnote{编者注: 从\SGARefLink{sga3:iii:paragraph:4-16}{4.16} 起, 情形会简化:
我们将只考虑在 \(S\) 上平坦的概形, \(Y\) 将是平坦的
\(S\)-群, 且 \(Y'=Y\times_S Y\); 这将给出 \(S_0\)-函子
\(N_0\) 和 \(N'_0\).}

% Source page: Exp3 local p. 57; combined p. 167; original printed p. 138.
在任何情形下都有交换图

% Raster stand-in for Loop 1; replace by native TeX/TikZ in Loop 2.
% Authority: Exp3-14oct24.pdf local p. 57; combined reader p. 167.
\begin{center}
  \SGAThreeNativeInput[0.65\textwidth]
    {figures/exp3/diagram_045_exp3_p057_412_tensor_square.tex}
\end{center}

由于 \(Y'\) 在 \(S\) 上平坦, 左侧竖直箭头是同构,
从而得到 \(\mathscr O_{T_0}\)-模态射
\[
 f_0^*(\mathscr J\mathscr O_{Y'})\longrightarrow
 \mathscr J\mathscr O_T.
\]
它诱导交换群态射
\begin{multline*}
 \operatorname{Hom}_{\mathscr O_{T_0}}
 \bigl(f_0^*g_0^*(\mathscr N_{Y/X}
  \otimes_{\mathscr O_Y}\mathscr O_{Y_0}),
  f_0^*(\mathscr J\mathscr O_{Y'})\bigr)\\
 \longrightarrow M(T),
\end{multline*}
将它与下列由 \(f_0^*\) 诱导的态射
\begin{multline*}
 M(Y')\longrightarrow
 \operatorname{Hom}_{\mathscr O_{T_0}}
 \bigl(f_0^*g_0^*(\mathscr N_{Y/X}
  \otimes_{\mathscr O_Y}\mathscr O_{Y_0}),
  f_0^*(\mathscr J\mathscr O_{Y'})\bigr)
\end{multline*}
复合, 就得到态射 \(b_{f_0}:M(Y')\to M(T)\).

同样, 在任何情形下都有图

% Raster stand-in for Loop 1; replace by native TeX/TikZ in Loop 2.
% Authority: Exp3-14oct24.pdf local p. 57; combined reader p. 167.
\begin{center}
  \SGAThreeNativeInput[0.67\textwidth]
    {figures/exp3/diagram_046_exp3_p057_412_conormal_square.tex}
\end{center}

由于 \(Y'\) 在 \(S\) 上平坦, 由\SGARefLink{sga3:iii:remark:4-8-0:equation:4-8-0}{4.8.0},
第二个竖直箭头为同构. 因此得到 \(\mathscr O_{Y'_0}\)-态射
\[
 i_0'^*g_0^*(\mathscr N_{Y/X}
   \otimes_{\mathscr O_Y}\mathscr O_{Y_0})
 \longrightarrow
 \mathscr N_{Y'/X'}
   \otimes_{\mathscr O_{Y'}}\mathscr O_{Y'_0}.
\]
这诱导态射 \(a_{g_0}(\operatorname{id}_{Y'}):M'(Y')\to M(Y')\),
且对每个 \(f:T\to Z\), 诱导态射 \(a_{g_0}(f):M'(T)\to M(T)\),
使下图交换:

% Raster stand-in for Loop 1; replace by native TeX/TikZ in Loop 2.
% Authority: Exp3-14oct24.pdf local p. 57; combined reader p. 167.
\begin{center}
  \SGAThreeNativeInput[0.50\textwidth]
    {figures/exp3/diagram_047_exp3_p057_412_functor_square.tex}
\end{center}

这里 \(b'_{f_0}\) 的定义与 \(b_{f_0}\) 相同.
\hfill\textsc{证毕.}

\medskip
\noindent\SGAInlineRefTarget{sga3:iii:remark:4-12-1}\textbf{注 4.12.1.}
\footnote{编者注: 此注为新增内容, 将用于\SGARefLink{sga3:iii:lemma:4-17}{4.17} 的证明.}
\textit{设 \(M_0\) 和 \(M'_0\) 为如下定义的 \(Y'_0\)-函子:
对每个 \(f:T_0\to Y'_0\), 令}
\begin{align*}
 M_0(T)&=
 \operatorname{Hom}_{\mathscr O_{T_0}}
 \bigl(f_0^*g_0^*(\mathscr N_{Y/X}
  \otimes_{\mathscr O_Y}\mathscr O_{Y_0}),
  \mathscr J\otimes_{\mathscr O_{S_0}}\mathscr O_{T_0}\bigr),\\
 M'_0(T)&=
 \operatorname{Hom}_{\mathscr O_{T_0}}
 \bigl(f_0^*(\mathscr N_{Y'/X'}
  \otimes_{\mathscr O_{Y'}}\mathscr O_{Y'_0}),
 \mathscr J\otimes_{\mathscr O_{S_0}}\mathscr O_{T_0}\bigr).
\end{align*}

% Source page: Exp3 local p. 58; combined p. 168; original printed p. 139.
立即可见 \(Z_0=Y'_0\), 且在 \(S\) 上平坦的
\(Z\)-概形 \(T\) 所组成的范畴上, 函子 \(M\) 和 \(M'\)
分别与 \(\prod_{S_0/S}M_0\) 和 \(\prod_{S_0/S}M'_0\) 一致.
此时, \(b_{f_0}\) 就是由 \(f_0\) 诱导的态射
\[
 f_0^*:M_0(Y'_0)\longrightarrow M(T_0)
\]
此外, 对每个态射 \(u:U\to T\),
若 \(h=f\circ u\), 则下图交换:

% Raster stand-in for Loop 1; replace by native TeX/TikZ in Loop 2.
% Authority: Exp3-14oct24.pdf local p. 58; combined reader p. 168.
\begin{center}
  \SGAThreeNativeInput[0.50\textwidth]
    {figures/exp3/diagram_048_exp3_p058_4121_naturality_square.tex}
\end{center}

因此 \(a_{g_0}\) 成为函子态射
\[
 \prod_{S_0/S}M'_0\longrightarrow\prod_{S_0/S}M_0.
\]

\medskip
\noindent\SGAInlineRefTarget{sga3:iii:proposition:4-13}\textbf{命题 4.13.}
\textit{假设 \(Y'\) 在 \(S\) 上平坦. 则}
\[
 c(X,Y,g\circ f)=a_{g_0}\bigl(c(X',Y',f)\bigr)
   +b_{f_0}\bigl(c(X,Y,g\circ i')\bigr).
\]

由于各个障碍以及态射 \(a_{g_0}\) 和 \(b_{f_0}\) 的定义都是局部的,
显然只需在所涉及的概形全为仿射概形时验证此公式. 记
\[
 \begin{gathered}
 S=\operatorname{Spec}(\Lambda),\qquad
 S_{\mathscr J}=\operatorname{Spec}(\Lambda/\mathscr J),\qquad
 S_0=\operatorname{Spec}(\Lambda/\mathscr I),\\
 T=\operatorname{Spec}(C),\qquad X=\operatorname{Spec}(A),\qquad
 Y=\operatorname{Spec}(A/I_Y)=\operatorname{Spec}(B),\\
 X'=\operatorname{Spec}(A'),\qquad
 Y'=\operatorname{Spec}(A'/I_{Y'})=\operatorname{Spec}(B').
 \end{gathered}
\]
因此有如下环与理想的图
\footnote{编者注: 沿用原版的记号, 用 \(f:A'\to C\) 和 \(g:A\to A'\)
表示对应于 \(f:T\to X'\) 和 \(g:X'\to X\) 的环同态.
这就解释了对 \(x\in I_Y\) 的公式
\(c(X,Y,g\circ f)(\bar x\otimes1)=f(g(x))\).}

% Raster stand-in for Loop 1; replace by native TeX/TikZ in Loop 2.
% Authority: Exp3-14oct24.pdf local p. 58; combined reader p. 168.
\begin{center}
  \SGAThreeNativeInput[0.48\textwidth]
    {figures/exp3/diagram_049_exp3_p058_413_ring_ideal_diagram.tex}
\end{center}

考察待证公式中的各项. 下文中, 对 \(x\in I_Y\)
(相应地, \(u\in I_{Y'}\)), 用 \(\bar x\) (相应地, \(\bar u\))
表示它在 \(I_Y/I_Y^2\) (相应地, \(I_{Y'}/I_{Y'}^2\)) 中的像.
若 \(m\) 属于 \(\Lambda\)-模 \(M\), 则用 \(\bar m_0\)
表示其在 \(M_0=M/\mathscr I M\) 中的像.

我们已知 \(c=c(X,Y,g\circ f)\) 是满足随后条件的唯一 \(C_0\)-态射
\[
 I_Y/I_Y^2\otimes_B C_0\longrightarrow\mathscr J C
\]
该条件为对每个 \(x\in I_Y\), 均有 \(c(\bar x\otimes1)=f(g(x))\).

% Source page: Exp3 local p. 59; combined p. 169; original printed p. 140.
固定 \(x\in I_Y\). 由于 \(g_{\mathscr J}(Y'_{\mathscr J})\subset
Y_{\mathscr J}\), 有 \(g(x)\in I_{Y'}+\mathscr J A'\). 记
\[
 g(x)=x'+\sum_i\lambda_i a'_i,
 \qquad x'\in I_{Y'},\quad \lambda_i\in\mathscr J,\quad a'_i\in A'.
\]
于是
\begin{equation*}
 \tag{1}\label{sga3:iii:proposition:4-13:equation:1}
 c(X,Y,g\circ f)(\bar x\otimes1)
 =f(g(x))=f(x')+\sum_i\lambda_i f(a'_i).
\end{equation*}

考虑 \(a_{g_0}(c(X',Y',f))\). 根据刚给出的定义, 它由下图定义:

% Raster stand-in for Loop 1; replace by native TeX/TikZ in Loop 2.
% Authority: Exp3-14oct24.pdf local p. 59; combined reader p. 169.
\begin{center}
  \SGAThreeNativeInput[0.86\textwidth]
    {figures/exp3/diagram_050_exp3_p059_413_ag0_obstruction_diagram.tex}
\end{center}

因此 \(a_{g_0}(c(X',Y',f))(\bar x\otimes1)=f(u)\),
其中 \(u\in I_{Y'}\) 在 \(I_{Y'_0}/I_{Y'_0}^2\) 中的像 \(\bar u\) 满足
\[
 \bar u_0\otimes1=\bar g_0(\bar x_0)\otimes1
 =g_0(\bar x_0)\otimes1.
\]
故可取 \(u=x'\), 得到
\begin{equation*}
 \tag{2}\label{sga3:iii:proposition:4-13:equation:2}
 a_{g_0}(c(X',Y',f))(\bar x\otimes1)=f(x').
\end{equation*}

最后考虑 \(b_{f_0}(c(X,Y,g\circ i'))\). 根据假设,
\(\Lambda_0\)-代数态射 \(f_0:A'_0\to C_0\) 经由 \(B'_0\) 分解.
由于 \(Y'\) 平坦, 因而有
\[
 \mathscr J\otimes_{\Lambda_0}B'_0\xrightarrow{\sim}\mathscr J B',
\]
于是得到 \(B'_0\)-模态射 \(\psi:\mathscr J B'\to\mathscr J C\),
使下图交换:

% Raster stand-in for Loop 1; replace by native TeX/TikZ in Loop 2.
% Authority: Exp3-14oct24.pdf local p. 59; combined reader p. 169.
\begin{center}
  \SGAThreeNativeInput[0.74\textwidth]
    {figures/exp3/diagram_051_exp3_p059_413_tensor_J_square.tex}
\end{center}

设
\[
 \varphi:\mathscr J B'\otimes_{B'_0}C_0\longrightarrow\mathscr J C
\]
为 \(\psi\) 诱导的 \(C_0\)-模态射. 则对每个
\(a'\in A'\) 及 \(\lambda\in\mathscr J\),
\begin{equation*}
\tag{\(\dagger\)}\label{sga3:iii:proposition:4-13:equation:dagger}
 \varphi(\lambda\pi'(a')\otimes1)=\lambda f(a').
\end{equation*}

% Source page: Exp3 local p. 60; combined p. 170; original printed p. 141.
于是元素 \(b_{f_0}(c(X,Y,g\circ i'))\) 由如下交换图定义:

% Raster stand-in for Loop 1; replace by native TeX/TikZ in Loop 2.
% Authority: Exp3-14oct24.pdf local p. 60; combined reader p. 170.
\begin{center}
  \SGAThreeNativeInput[0.62\textwidth]
    {figures/exp3/diagram_052_exp3_p060_413_bf0_obstruction_diagram.tex}
\end{center}

因此,
\begin{equation*}
 \tag{3}\label{sga3:iii:proposition:4-13:equation:3}
 b_{f_0}(c(X,Y,g\circ i'))(\bar x\otimes1)
 =\varphi\left(\sum_i\lambda_i\pi'(a'_i)\otimes1\right)
 =\sum_i\lambda_i f(a'_i),
\end{equation*}
最后一个等式由 \((\dagger)\) 得到.
比较三个显式公式\SGARefLink{sga3:iii:proposition:4-13:equation:1}{(1)},
\SGARefLink{sga3:iii:proposition:4-13:equation:2}{(2)} 和\SGARefLink{sga3:iii:proposition:4-13:equation:3}{(3)},
即得所求公式.

\medskip
\noindent\SGAInlineRefTarget{sga3:iii:corollary:4-14}\textbf{推论 4.14.}
\textit{设 \(Y,Y'\) 为 \(X\) 的两个平坦子概形,
均约化为 \(Y_{\mathscr J}\), 并假设 \(Y_0\) 在 \(X_0\) 中是局部完全交.
若 \(f:T\to X\) 为 \(S\)-态射, 使得
\(f_{\mathscr J}\) 经由 \(Y_{\mathscr J}\to X_{\mathscr J}\) 分解, 则}
\[
 c(X,Y,f)-c(X,Y',f)=b_{f_0}\bigl(d(Y,Y')\bigr).
\]

事实上, 将上述公式应用于图

% Raster stand-in for Loop 1; replace by native TeX/TikZ in Loop 2.
% Authority: Exp3-14oct24.pdf local p. 60; combined reader p. 170.
\begin{center}
  \SGAThreeNativeInput[0.62\textwidth]
    {figures/exp3/diagram_053_exp3_p060_414_comparison_diagram.tex}
\end{center}

得到 \(c(X,Y,f)-c(X,Y',f)=b_{f_0}(c(X,Y,i'))\).
此外, 由\SGARefLink{sga3:iii:proposition:4-8:item:ii}{4.8(ii)},
\(c(X,Y,i')=d(Y,Y')\).

\medskip
\noindent\SGAInlineRefTarget{sga3:iii:proposition:4-15}\textbf{命题 4.15.}
\textit{设 \(X\) 为在 \(S\) 上平滑的 \(S\)-群,
\(Y\) 为其在 \(S\) 上平坦且局部有限呈示的子 \(S\)-群.
则 \(Y\) 在 \(X\) 中是局部完全交 (参见\SGARefLink{sga3:iii:definition:4-6-2}{4.6.2}).}

% Source page: Exp3 local p. 61; combined p. 171; original printed p. 142.
\noindent\textit{证明.}
\footnote{编者注: 在陈述中加入了 \(Y\) 在 \(S\) 上局部有限呈示的假设,
并补出了下面的证明, 它比原版所勾画的证明更直接.
为完整起见, 也将后者写出. 与上面的证明一样, 首先归结为
\(S=\operatorname{Spec}(k)\), 其中 \(k\) 代数闭的情形.
由 EGA IV\textsubscript{4}, 16.9.10 和 19.3.2,
只需证明对每个 \(y\in Y\), 局部环 \(\mathscr O_{Y,y}\) 的完备化是一个完备
Noether 局部环对某个正则序列的商. 由同上 19.3.3,
满足此性质的 \(y\in Y\) 所组成的集合是开子集 \(U\subset Y\).
由于 \(Y\) 在 \(k\) 上有限型, 只需证明 \(U\) 包含每个闭点.
由于 \(Y\) 为 \(k\)-群, 经平移可归结为考察
\(\mathscr O_{Y,e}\) 的完备化, 即 \(Y\) 对应的形式群 \(\widehat Y\)
(参见第 VII\textsubscript{B} 篇). 由于 \(X\) 平滑,
仿射代数 \(\mathscr A(\widehat X)\) 是幂级数代数 \(k[[X_1,\ldots,X_n]]\).
Dieudonn\'e 结构定理表明 \(\mathscr A(\widehat Y)\) 同构于商
\[
 k[[X_1,\ldots,X_{r+s}]]/
 (X_1^{p^{n_1}},\ldots,X_r^{p^{n_r}}),
\]
其中 \(p\) 为 \(k\) 的特征指数, 且 \(r+s\leq n\);
参见第 VII\textsubscript{B} 篇, 注 5.5.2(b).}
我们证明浸入 \(Y\to X\) 在 EGA IV\textsubscript{4}, 16.9.2 的意义下正则.
由 EGA IV\textsubscript{4}, 19.5.1, 这蕴涵它在 4.6.2 的意义下也正则
(事实上, 由同上可知, 当 \(S\) 局部 Noether 时, 两个定义等价).
因此这里的\lq\lq 正则浸入\rq\rq 将指 EGA IV\textsubscript{4}, 16.9.2 意义下的正则浸入.
由于 \(X\) 和 \(Y\) 在 \(S\) 上平坦且局部有限呈示,
EGA IV\textsubscript{4}, 19.2.4 将此断言归结为证明:
对每个 \(s\in S\), \(Y_s\to X_s\) 为正则浸入.
由 EGA IV\textsubscript{4}, 19.1.5(ii),
只需在 \(S\) 的几何纤维上验证, 因而可假设 \(S\) 是代数闭域 \(k\) 的谱.

由第 VI\textsubscript{A} 篇\SGARefLink{thm:sga3-VIA-3.2}{3.2},
商 \(X/Y\) 存在且平滑, 态射 \(\pi:X\to X/Y\) 平坦, 且下图为笛卡尔方图:

% Raster stand-in for Loop 1; replace by native TeX/TikZ in Loop 2.
% Authority: Exp3-14oct24.pdf local p. 61; combined reader p. 171.
\begin{center}
  \SGAThreeNativeInput[0.35\textwidth]
    {figures/exp3/diagram_054_exp3_p061_415_cartesian_square.tex}
\end{center}

这里 \(\bar e\) 为 \(X\) 的单位点在 \(X/Y\) 中的像.
由平坦基变换 (参见 EGA IV\textsubscript{4}, 19.1.5(iii)),
只需证明 \(i\) 为正则浸入. 这显然成立:
Noether 局部环 \(\mathscr O_{X/Y,\bar e}\) 平滑 (从而正则),
所以其极大理想由正则序列生成.

\medskip
\noindent\SGAInlineRefTarget{sga3:iii:paragraph:4-16}\textbf{4.16.}
\footnote{编者注: 我们重新组织了\SGARefLink{sga3:iii:paragraph:4-16}{4.16},
将\SGARefLink{sga3:iii:proposition:4-15}{4.15} 末尾陈述的假设和障碍 \(DY\) 的定义集中在此.}
设 \(X\) 为在 \(S\) 上平滑的 \(S\)-群,
用 \(\mu:X\times_S X\to X\) 表示其群运算.
设 \(Y_{\mathscr J}\) 为 \(X_{\mathscr J}\) 的子 \(S_{\mathscr J}\)-群,
在 \(S_{\mathscr J}\) 上平坦且局部有限呈示.
由\SGARefLink{sga3:iii:proposition:4-15}{4.15},
\(Y_{\mathscr J}\) 在 \(X_{\mathscr J}\) 中是局部完全交.
从而由\SGARefLink{sga3:iii:proposition:4-6-5}{4.6.5}, 每个平坦的
\footnote{编者注: 已对原文作了修正, 加入\lq\lq 平坦\rq\rq 一词.}
提升 \(Y_{\mathscr J}\) 的 \(S\)-概形 \(Y\) 在 \(X\) 中都是局部完全交.
对这样的 \(Y\), 由\SGARefLink{sga3:iii:remark:4-8-0}{4.8.0} 得
\begin{equation*}
 \tag{4.16.1}\label{sga3:iii:paragraph:4-16:equation:4-16-1}
 \mathscr N_{Y/X}\otimes_{\mathscr O_Y}\mathscr O_{Y_0}
 =\mathscr N_{Y_0/X_0}
 =\mathscr N_{Y_{\mathscr J}/X_{\mathscr J}}
   \otimes_{\mathscr O_{Y_{\mathscr J}}}\mathscr O_{Y_0}.
\end{equation*}

% Source page: Exp3 local p. 62; combined p. 172; original printed p. 142.
另一方面, 设 \(\varepsilon_0:S_0\to Y_0\) 为 \(Y_0\) 的单位截面,
并设 \(\mathbf n_{Y_0/X_0}\) 为拟凝聚 \(\mathscr O_{S_0}\)-模
\[
\mathbf n_{Y_0/X_0}=\varepsilon_0^*(\mathscr N_{Y_0/X_0}).
\]
由于 \(Y_0\) 和 \(X_0\) 都是 \(S_0\)-群, 显然
\(\mathscr N_{Y_0/X_0}\) 在例如 \(Y_0\) 的左平移下不变. 因此,
\footnote{编者注: 见下文\SGARefLink{sga3:iii:lemma:4-25}{4.25}.}
它是 \(\mathbf n_{Y_0/X_0}\) 在 \(Y_0\to S_0\) 下的逆像; 即
\begin{equation*}
 \tag{4.16.2}\label{sga3:iii:paragraph:4-16:equation:4-16-2}
 \mathscr N_{Y_0/X_0}
=\mathbf n_{Y_0/X_0}\otimes_{\mathscr O_{S_0}}\mathscr O_{Y_0}.
\end{equation*}

结合 (\SGARefLink{sga3:iii:paragraph:4-16:equation:4-16-1}{4.16.1})
和 (\SGARefLink{sga3:iii:paragraph:4-16:equation:4-16-2}{4.16.2}),
首先由\SGARefLink{sga3:iii:proposition:4-5}{4.5} 推出,
\(X\) 中在 \(S\) 上平坦且提升 \(Y_{\mathscr J}\) 的子概形 \(Y\) 所组成的集合,
或为空集, 或为下列群下的主齐性空间:
\begin{equation*}
 \tag{4.16.3}\label{sga3:iii:paragraph:4-16:equation:4-16-3}
 \operatorname{Hom}_{\mathscr O_{Y_0}}
\bigl(\mathbf n_{Y_0/X_0}\otimes_{\mathscr O_{S_0}}\mathscr O_{Y_0},
       \mathscr J\otimes_{\mathscr O_{S_0}}\mathscr O_{Y_0}\bigr).
\end{equation*}
另一方面, \SGARefLink{sga3:iii:proposition:4-8:item:i}{4.8(i)} 表明,
对每个这样的 \(Y\) 及每个满足
\(f_{\mathscr J}:T_{\mathscr J}\to X_{\mathscr J}\) 经由
\(Y_{\mathscr J}\) 分解的 \(S\)-态射 \(f:T\to X\),
\(f\) 经由 \(Y\) 分解的障碍 \(c(X,Y,f)\) 属于
\[
 \operatorname{Hom}_{\mathscr O_{T_0}}
\bigl(\mathbf n_{Y_0/X_0}\otimes_{\mathscr O_{S_0}}\mathscr O_{T_0},
       \mathscr J\mathscr O_T\bigr).
\]
若 \(T\) 也在 \(S\) 上平坦, 则最后这个群等于
\[
 \operatorname{Hom}_{\mathscr O_{T_0}}
\bigl(\mathbf n_{Y_0/X_0}\otimes_{\mathscr O_{S_0}}\mathscr O_{T_0},
       \mathscr J\otimes_{\mathscr O_{S_0}}\mathscr O_{T_0}\bigr).
\]
这促使我们引入下述群函子 \(N_0\).

\medskip
\noindent\SGAInlineRefTarget{sga3:iii:definition:4-16-1}\textbf{定义 4.16.1.}
\textit{设 \(N_0\) 为如下定义的交换群值 \(S_0\)-函子:
对每个 \(Z\in\operatorname{Ob}(\mathbf{Sch})_{/S_0}\), 令}
\begin{equation*}
 \tag{*}\label{sga3:iii:definition:4-16-1:equation:star}
 \operatorname{Hom}_{S_0}(Z,N_0)
 =\operatorname{Hom}_{\mathscr O_Z}
\bigl(\mathbf n_{Y_0/X_0}\otimes_{\mathscr O_{S_0}}\mathscr O_Z,
       \mathscr J\otimes_{\mathscr O_{S_0}}\mathscr O_Z\bigr).
\end{equation*}
\textit{则 \(X\) 中在 \(S\) 上平坦且提升 \(Y_{\mathscr J}\) 的子概形
\(Y\) 所组成的集合, 或为空集, 或为下列群下的主齐性空间:}
\[
 \operatorname{Hom}_{S_0}(Y_0,N_0)=C^1(Y_0,N_0).
\]

对每个这样的 \(Y\), 考虑图

% Raster stand-in for Loop 1; replace by native TeX/TikZ in Loop 2.
% Authority: Exp3-14oct24.pdf local p. 62; combined reader p. 172.
\begin{center}
  \SGAThreeNativeInput[0.52\textwidth]
    {figures/exp3/diagram_055_exp3_p062_4161_DY_diagram.tex}
\end{center}

并令
\[
 DY=c\bigl(X,Y,\mu\circ(i,i)\bigr)
\]
为 \(\mu\circ(i,i)\) 经由 \(Y\) 分解的障碍,
即 \(Y\) 在 \(X\) 的群运算下稳定的障碍.
由上文, \(DY\) 是下列群的元素:
\[
 N_0(Y_0\times_{S_0}Y_0)=C^2(Y_0,N_0).
\]

% Translation coverage: entire assigned source component, 577 source lines.
% Formula audit: 277 expressions; math unchanged except translated text annotations.
% All source reference keys, labels, footnotes and native diagram calls retained.
